\section{Vector V1: Parameter Manipulation}
\label{sec:v1}

\subsection{Attack Description}

The \bisect{} algorithm exposes five tunable parameters: METIS imbalance
factor $u_{\rm factor}$, county edge weight $\alpha_{\rm county}$,
\cs{} threshold $T$, area swing $a_{\rm swing}$, and VRA boost weight
$w_{\rm vra}$.
A hostile actor could observe partisan implications of each parameter setting
and then select parameter values that favor their preferred party while
claiming the choices are technical rather than political.

The attack is particularly insidious because each parameter has a legitimate
technical justification:
$u_{\rm factor}$ controls population balance tolerance;
$\alpha_{\rm county}$ controls county integrity preservation;
$T$ controls convergence completeness;
$a_{\rm swing}$ controls area-vs-compactness tradeoff in GeoSection;
$w_{\rm vra}$ controls minority community protection weighting.

\subsection{Maximum Achievable Shift}

U.2's 50-state sweep~\citep{b17paper} provides the definitive empirical
answer.
Across all five parameters varied independently from most Democrat-favorable
to most Republican-favorable settings:
\begin{itemize}
\item $u_{\rm factor}$: $\Delta D \leq 0.4$ nationally (0.1\% vs.\ 5\%)
\item $\alpha_{\rm county}$: $\Delta D \leq 0.2$ nationally (1.0 vs.\ 10.0)
\item $T$: $\Delta D \leq 0.1$ nationally (100 vs.\ 1000)
\item $a_{\rm swing}$: $\Delta D \leq 0.2$ nationally (1.05 vs.\ 1.20)
\item $w_{\rm vra}$: $\Delta D \leq 0.2$ nationally (0.2 vs.\ 0.6)
\end{itemize}
The combined maximum (additive bound): $\Delta D \leq 0.5$ seats.
The empirical combined effect from U.2's joint sweep: $\Delta D \leq 0.3$ seats.

This is comparable to the seed-variance noise floor from B.7, which
finds a coefficient of variation below 2\% in partisan outcomes across
random seed choices.

\subsection{Why the Effect Is So Small}

The insensitivity of partisan outcomes to parameter tuning reflects the
following structural features of the algorithm (U.2, Section~5):
\begin{enumerate}
\item \textbf{Population neutrality}: Parameters affect how well population
  constraints are satisfied, not which tracts land in which districts.
\item \textbf{Geographic continuity}: The bisection tree structure is
  determined by the prime factorisation of the district count $k$ and the
  graph topology, not by parameter values.
  Parameter changes alter leaf-level cut quality within a fixed tree,
  not the tree itself.
\item \textbf{Competitive-district geography}: Partisan outcomes depend
  primarily on how the algorithm handles geographically marginal competitive
  regions.
  A parameter change that uniformly shifts district boundaries by 1--2
  tracts will not systematically favor one party in swing regions.
\end{enumerate}

\subsection{Detection Mechanism}

Under the \dia{} protocol, parameters are fixed in a publicly registered
formula before Census redistricting data are released.
Any parameter deviation from the formula is detectable by comparing
the deployment manifest against the registered formula.
The manifest includes a SHA-256 hash of the parameter configuration file;
any change to any parameter changes the hash.
Detection probability: $\delta_1 = 1.00$ (cryptographic guarantee).
